\section{Background: Interval Type-2 Fuzzy Systems}
\label{sec:background}

A Type-1 fuzzy set assigns each input $x$ a crisp membership
$\mu(x)\in[0,1]$~\cite{zadeh1975linguistic}. When the shape or position of
that set is itself uncertain, a Type-2 set assigns a fuzzy membership
instead. In the interval case, every primary membership is an interval
$[\mulow(x),\muup(x)]$, and the region between the lower and upper membership
functions is the footprint of uncertainty
(FOU)~\cite{liang2000it2,mendel2006it2simple}.

A common way to build an IT2 set is a Gaussian with an uncertain mean
$m\in[m_1,m_2]$ and fixed spread $s$:
\begin{equation}
\muup(x)=
\begin{cases}
\mathcal{N}(x;m_1,s) & x<m_1\\
1 & m_1\le x\le m_2\\
\mathcal{N}(x;m_2,s) & x>m_2
\end{cases}
\qquad
\mulow(x)=\min\{\mathcal{N}(x;m_1,s),\,\mathcal{N}(x;m_2,s)\},
\label{eq:um-gauss}
\end{equation}
where $\mathcal{N}(x;m,s)=\exp\!\big(-\tfrac12((x-m)/s)^2\big)$. The width
$m_2-m_1$ controls how much the system hedges; when it is zero the set
collapses to a Type-1 Gaussian.

For a rule base with consequents $y_i$, each rule $i$ fires with an interval
$[\underline{f}_i,\overline{f}_i]$ obtained by applying a t-norm to the lower
and to the upper memberships of its antecedents. Type reduction computes the
centroid interval
\begin{equation}
y_l=\min_{f_i\in[\underline{f}_i,\overline{f}_i]}
  \frac{\sum_i f_i y_i}{\sum_i f_i},\qquad
y_r=\max_{f_i\in[\underline{f}_i,\overline{f}_i]}
  \frac{\sum_i f_i y_i}{\sum_i f_i},
\label{eq:km}
\end{equation}
which the Karnik--Mendel (KM) procedure finds exactly by locating one switch
point in the sorted consequents~\cite{karnik2001centroid}. A crisp output is
usually $(y_l+y_r)/2$, but the interval itself is informative: its width
reflects how much the input uncertainty affects the conclusion~\cite{wu2002bounds}.
IT2 controllers have been used in mobile robotics to cope with noisy
sensors and changing environments~\cite{hagras2004hierarchical}, typically
with an FOU chosen offline. In this work the FOU width is an online signal.
